%!TEX root = ../../../thesis-trier.tex

Here, in \cref{thm:hom_lower_bound_complexity},
we show that transforming a constraint language
yields a problem whose complexity is
at most that of the original problem w.r.t.~$\leACZEROm$-reductions.
In \cref{thm:constr_sync_hom_strongly_self_sync}
we show that for coding morphisms of solid codes
the complexity is preserved.

\begin{theorem}
	\label{thm:hom_lower_bound_complexity}
	Let $\varphi \colon \Sigma^* \to \Gamma^*$ be a fixed morphism
	and $L \subseteq \Sigma^*$.
	Then \[ \constrsync{\varphi(L)} \le_m^{\AC^0} \constrsync{L}. \]
\end{theorem}
\begin{proof}
	For the $\leACZEROm$-reduction,
	let $\semiautomat{A}{Q}{\Gamma}{\delta}$ be an input automaton. 
	We want to know if $\mathcal A$ is synchronizing with respect to~$\varphi(L)$.
	
	\medskip
	\noindent\underline{Construction:} Let $\semiautomat{A'}{Q}{\Sigma}{\delta'}$
	with $\applytrans{\delta'}{q}{x} = \applytrans{\delta}{q}{\applyfunc{x}{\varphi}}$ for $q \in Q$ and $x \in \Sigma$.
	We obtain a representation of $\mathcal A'$
	by transition matrices by multiplying
	for each $x \in \Sigma$
	the respective transition matrices
	of the letters in~$\applyfunc{x}{\varphi}$,
	which is doable by a \LOGSPACE-uniform circuit family in $\AC^0$ as $\varphi$ is fixed and matrix multiplication for a fixed number of matrices is in $\AC^0$ by \cref{prelim:circuit:ex:AC0mtx}.
		
	\medskip 
	\noindent\underline{Verification:} We have
	\begin{align*}
		                  \exists u \in \varphi(L) \colon |\delta(Q, u)| = 1 
		& \Leftrightarrow \exists v \in L \colon |\delta(Q, \varphi(v))| = 1 \tag{$\varphi(v) = u$} \\
		& \Leftrightarrow \exists v \in L \colon |\delta'(Q, v)| = 1  \tag{$\delta'(Q, v) = \delta(Q, \varphi(v))$}.  
	\end{align*}
So $\mathcal A'$ has a synchronizing word in $L$ if and only if
$\mathcal A$ has a synchronizing word in $\applyfunc{L}{\varphi}$.
\end{proof}

\begin{example}
	Let $L = \{a,b\}^* \cup \{a,c\}^* \cup \{b,c\}^*$.
	Then $\pi_{\{a,b\}}(L) = \{a,b\}^*$,
	so that $\constrsync{L}$ is \NL-hard
	with \cref{thm:hom_lower_bound_complexity}
	and \crefitem{constr_sync:thm:introductory_examples}{constrsync:ex:item:ab}.
	With \cref{constr_sync:thm:union} 
	we also get membership in \NL.
\end{example}	

For coding morphisms of solid codes,
we have a stronger statement than 
Theorem~\ref{thm:hom_lower_bound_complexity}.

\begin{theorem}
	\label{thm:constr_sync_hom_strongly_self_sync}
	Let $\varphi \colon \Sigma^* \to \Gamma^*$
	be a coding morphism such that\footnote{Recall that $\Sigma$ is finite, and so $C$ is finite. Also coding morphism
		are injective by definition.}
	$C = \varphi(\Sigma)$ 
	is a solid code.
	Then, for each regular $L \subseteq \Sigma^*$ we have 
	$
	\constrsync{L} \eqACZEROm \constrsync{\varphi(L)}.
	$
\end{theorem}
\begin{proof}
	By Theorem~\ref{thm:hom_lower_bound_complexity}, we have 
	$\constrsync{\varphi(L)} \leACZEROm \constrsync{L}$.
	Next, we show $\constrsync{L} \leACZEROm \constrsync{\applyfunc{L}{\varphi}}$
	by utilizing the Levenshtein-Romanov mappping from \cref{constr_sync:subsec:levenshtein-romanov-map}
    w.r.t.~the language~$C$.
	Now, let $\semiautomat{A}{Q}{\Sigma}{\delta}$ be an input automaton.
	Suppose $\Sigma = \{a_1, \ldots, a_k\}$
	and set $u_i = \applyfunc{a_i}{\varphi}$ for $i \in \{1,\ldots,k\}$.
	Then $C = \{u_1, \ldots, u_k\}$.	
	
	\begin{figure}[htbp]
		\centering		
		\begin{subfigure}{0.33\textwidth}
			\centering
	\begin{tikzpicture}[>=stealth,shorten >=1pt,auto]
	\node[state] (q0) at (90:2.5cm) {$q_0$};
	\node[state] (q1) at (210:2.5cm) {$q_1$};
	\node[state] (q2) at (330:2.5cm) {$q_2$};
	
	\path[->]
	(q0) edge[bend left=10] node[right] {$a_1$} (q2)
	(q2) edge[bend left=10] node[below] {$a_1$} (q1)
	(q1) edge[bend left=5]  node[left] {$a_1$}  (q0);
	
	\path[->] (q1) edge[bend left=10] node {$a_2$} (q2)
	(q2) edge[bend left=10] node {$a_2$} (q0)
	(q0) edge [loop above] node {$a_2$} (q0);
\end{tikzpicture}
			\caption{A base automaton~$\mathcal A$}
		\end{subfigure}
		\hfill
		\begin{subfigure}{0.55\textwidth}
		\centering
	\begin{tikzpicture}[>=stealth,shorten >=1pt, node distance=0.97cm, auto]
	
	%\tikzset{state/.style={}}
	\tikzset{state/.style={inner sep=2pt}}
	
	%%%%%%%%%%%%%%%%%%%%%%%%%%%%%%%%%%%%
	% TOP LAYER (epsilon)
	%%%%%%%%%%%%%%%%%%%%%%%%%%%%%%%%%%%%
	\node[state] (q0e) at (90:2.7cm) {$(q_0,\varepsilon)$};
	\node[state] (q1e) at (210:2.7cm) {$(q_1,\varepsilon)$};
	\node[state] (q2e) at (330:2.7cm) {$(q_2,\varepsilon)$};
	
	%%%%%%%%%%%%%%%%%%%%%%%%%%%%%%%%%%%%
	% MIDDLE LAYER (a)
	%%%%%%%%%%%%%%%%%%%%%%%%%%%%%%%%%%%%
	\node[state] (q0a) [below=of q0e] {$(q_0,a)$};
	\node[state] (q1a) [below=of q1e] {$(q_1,a)$};
	\node[state] (q2a) [below=of q2e] {$(q_2,a)$};
	
	%%%%%%%%%%%%%%%%%%%%%%%%%%%%%%%%%%%%
	% LOWER LAYER (aa)
	%%%%%%%%%%%%%%%%%%%%%%%%%%%%%%%%%%%%
	\node[state] (q0aa) [below=of q0a] {$(q_0,aa)$};
	\node[state] (q1aa) [below=of q1a] {$(q_1,aa)$};
	\node[state] (q2aa) [below=of q2a] {$(q_2,aa)$};
	
	\path[->]
	(q0e) edge (q0a)
	(q0a) edge (q0aa)
	(q0aa) edge[in=230,out=260, looseness=9] (q0aa)
	(q0aa) edge[dotted, in=280, out=310, looseness=9] (q0aa)
	
	(q1e) edge  (q1a)
	(q1a) edge  (q1aa)
	(q1aa) edge[loop below] ()
	
	(q2e) edge  (q2a)
	(q2a) edge  (q2aa)
	(q2aa) edge[loop below] ();		
	
	\path[->]
	(q0e)  edge[dashed, in=80, out=50, looseness=9] (q0e)
	(q0e)  edge[dotted, in=130, out=100, looseness=9] (q0e)
	(q0a)  edge[dotted, bend left=20] (q0e)
	(q0aa) edge[dashed, bend right=10] (q2e);
	
	\path[->]
	(q1e)  edge[dashed, in=80, out=50, looseness=9] (q1e)
	(q1e)  edge[dotted, in=130, out=100, looseness=9] (q1e)
	(q1a)  edge[dotted, bend left=20] (q1e)
	(q1aa) edge[dashed, bend left=65] (q0e);
	
	\path[->]
	(q2e)  edge[dashed, in=80, out=50, looseness=9] (q2e)
	(q2e)  edge[dotted, in=130, out=100, looseness=9] (q2e)
	(q2a)  edge[dotted, bend right=20] (q2e)
	(q2aa) edge[dashed, bend left=30] (q1e);		
	
	\path[->] 
	(q0a) edge[dashed, bend right=10] (q2e)
	(q2a) edge[dashed,bend left=30] (q1e)
	(q1a) edge[dashed,bend left=55] (q0e);	
	
	\path[->]
	(q1aa) edge[dotted, bend right=10] (q2e)
	(q2aa) edge[dotted, bend right=50] (q0e);		
\end{tikzpicture}
			\caption{Constructed automaton as in the proof.}
		\end{subfigure}		
		\caption[Construction from proof of \cref{thm:constr_sync_hom_strongly_self_sync}.]{Construction from the proof of \cref{thm:constr_sync_hom_strongly_self_sync}
		    for the solid code $\{ab, aac\}$
			and base automaton shown on the left side. Solid, dashed, and dotted lines represent transitions
			labeled with $a$, $b$, and~$c$, respectively.}
	\label{constrsync:fig:construction-solid-code}
	\end{figure}	
	
	\medskip
	\noindent\underline{Construction:} Define $\semiautomat{A'}{Q \times (\pref{C} \setminus C)}{\Gamma}{\delta'}$
	where, for $(q, u) \in Q \times \pref{C} \setminus C$
	and $x \in \Sigma$,
	\begin{equation}\label{constr_sync:eqn:aut_solid_code}
		\applytrans{\delta'}{(q,u)}{x} = \begin{cases}
			(q, \sigma_C(ux)) & \text{if } \sigma_C(ux) \notin C, \\
			(\applytrans{\delta}{q}{a_i}, \varepsilon) & \text{if } \exists i \in \{1,\ldots,k\} : \sigma_C(ux) = u_i. 
		\end{cases}
	\end{equation}
	By encoding the proper prefixes in the second component
	we keep track of the prefix scanned in the input so far.	
	See \cref{constrsync:fig:construction-solid-code} for an example construction
	for the solid code $C = \{ab,aac\}$.
	Note that in the example $\delta'((q, a), b) = \delta'((q,aa), b) = (\delta(q, a_1), \varepsilon)$
	since $\sigma_C(ab) = \sigma_C(aab) = ab$.
		
	
	\noindent\underline{Properties:} We give the proofs for Properties~\ref{constr_sync:solid_code_morphism:apply_codeword_proper_prefix} and~\ref{constr_sync:solid_code_morphism:apply_codeword},
	Property~\ref{constr_sync:solid_code_morphism:wordsvarphi} follows from Property~\ref{constr_sync:solid_code_morphism:apply_codeword} by induction
	as $\applyfunc{v}{\varphi} \in C^+$ for $v \in \Sigma^+$.
	\begin{enumerate}
		\item \label{constr_sync:solid_code_morphism:apply_codeword_proper_prefix}
		For $(q, u) \in Q \times ( \pref{C} \setminus C )$ and $v \in \pref{C} \setminus C$
		we have $\sigma_C(uv) \in \pref{C} \setminus C$ and $\applytrans{\delta'}{(q,u)}{v}  = (q, \sigma_C(uv))$.		
		
		\begin{proofitem}{\ref{constr_sync:solid_code_morphism:apply_codeword_proper_prefix}} 
		We show the claim by induction on $v$.
		As $\sigma_C(u) = u$, the claim is true for $v = \varepsilon$.
		Now, suppose $v = wx$ for $w \in \Sigma^*$ and $x \in \Sigma$.		
		As $C$ is overlap-free, 
		$wx \in \pref{C} \setminus C$ implies 
		$\sigma_C(uwx) \notin C$ and so $\sigma_C(uwx) \in \pref{C} \setminus C$;
		this implication is noted in \crefitem{constrsync:prop:sigma-solid-code-char}{constrsync:prop:sigma-solid-code-char:2}.
		By~\crefitem{constr_sync:prop:sigma_L_properties}{constr_sync:sigme_L:item:concatenation}
		we have $\sigma_C(\sigma_C(uw)x) = \sigma_C(uwx)$.
		So, for $q \in Q$, we obtain		
		\begin{align*}
			& \delta'((q,u), wx) \\
			& = \delta'(\delta'((q,u), w), x) & \tag{Def.~of $\delta'$} \\
			& = \delta'((q, \sigma_C(uw)), x) & \tag{Induction hypothesis} \\
			& = (q, \sigma_C(\sigma_C(uw)x)). & \tag{Eqn.~\eqref{constr_sync:eqn:aut_solid_code} as $\sigma_C(\sigma_C(uw)x) \notin C$} 
		\end{align*}
		Hence, we get $ \delta((q,u), wx) = (q, \sigma_C(uwx))$
        and the claim follows by induction.
     	\end{proofitem}
		
		\item \label{constr_sync:solid_code_morphism:apply_codeword}
		For $(q, u) \in Q \times ( \pref{C} \setminus C )$ and $i \in \{1,\ldots,k\}$ we have
		\[
		\applytrans{\delta'}{(q,u)}{u_i} = (\applytrans{\delta}{q}{a_i}, \varepsilon) 
		\]
		
		\begin{proofitem}{\ref{constr_sync:solid_code_morphism:apply_codeword}}
		Let $u_i = vx$ for $v \in \Sigma^*$ and $x \in \Sigma$.
		As $C$ is a prefix code, we have $v \in \pref{C} \setminus C$.		
		By \crefitem{constr_sync:prop:sigma_L_properties}{constr_sync:sigme_L:item:concatenation}
		we have $\sigma_C(\sigma_C(uv)x) = \sigma_C(uvx)$.
		By \crefitem{constrsync:prop:sigma-solid-code-char}{constrsync:prop:sigma-solid-code-char:3}
		as $vx \in C$, we find $\sigma_C(uvx) = vx = u_i$. Then
		\begin{align*}
			\delta'((q, u), vx) & = \delta'(\delta'((q,u), v), x)  & \tag{Def.~of $\delta'$} \\
			& = \delta'((q, \sigma_C(uv)), x) & \tag{Property~\ref{constr_sync:solid_code_morphism:apply_codeword_proper_prefix}} \\
			& = (\applytrans{\delta}{q}{a_i}, \varepsilon). & \tag{Eqn.~\eqref{constr_sync:eqn:aut_solid_code} as $\sigma_C(\sigma_C(uv)x)  = u_i$}
		\end{align*}
	    So, $\applytrans{\delta'}{(q,u)}{u_i} = (\applytrans{\delta}{q}{a_i}, \varepsilon)$.
		\end{proofitem}
		
		\item  \label{constr_sync:solid_code_morphism:wordsvarphi} For every $v \in \Sigma^+$ and $(q,u) \in Q \times (\pref{C} \setminus C)$,
		\begin{equation}\label{constr_sync:eqn:solid_code3}
			\applytrans{\delta'}{(q,u)}{\applyfunc{v}{\varphi}} = (\applytrans{\delta}{q}{v}, \varepsilon).
		\end{equation}
	\end{enumerate}	
	
	\medskip 
	\noindent\underline{Verification:} That the reduction is correct follows by the following
	equivalences:
	\begin{align*}
		& \exists w \in \varphi(L) \colon |\applytrans{\delta'}{Q \times (\pref{C}\setminus C)}{w}|  = 1 \\
		& \Leftrightarrow \exists v \in L \colon |\applytrans{\delta'}{Q \times (\pref{C}\setminus C)}{\varphi(v)}|  = 1 \\
		& \Leftrightarrow \exists v \in L \colon |\applytrans{\delta}{Q}{v} \times \{\varepsilon\}| = 1 \tag{Eqn.~\eqref{constr_sync:eqn:solid_code3}} \\
		& \Leftrightarrow \exists v \in L \colon |\applytrans{\delta}{Q}{v}| = 1. \qedhere
	\end{align*}
\end{proof}

\begin{remark} \label{constrsync:rem:alphabet_solid_special_case}
	The construction from \cref{constrsync:rem:alphabet} is a special case 
	of \cref{thm:constr_sync_hom_strongly_self_sync}
	as for $\Gamma, \Sigma$ with $\Gamma \subseteq \Sigma$
	the inclusion map $i \colon \Gamma \to \Sigma$
	given by $i(x) = x$ for $x \in \Gamma$
	is a coding morphism where $i(\Gamma) \subseteq \Sigma$ is a solid code.
\end{remark} 

\begin{remark}
	The construction from the proof of \cref{thm:constr_sync_hom_strongly_self_sync}
	with the transition function given by \cref{constr_sync:eqn:aut_solid_code}
	can be applied to any prefix code~$C$
	and yields a well-defined automaton~$\semiautomat{A'}{Q \times \pref{C}\setminus C}{\Sigma}{\delta'}$.	
	Can the assumption for~$C$ to be a solid code
	be weakened? 
	The crucial property for the reduction
	to work is that $\delta'(Q \times \pref{C} \setminus C, u_i) = \delta(Q, a_i) \times \{\varepsilon\}$
	for $i \in \{1,\ldots,k\}$ and $C = \{u_1, \ldots, u_k\}$.
	However, this is not guaranteed to hold if~$C$
	is merely a comma-free code.
	For example, consider $C = \{ u_1, u_2 \}$ with $u_1 = aab$ and $u_2 = bba$,
	which is a comma-free code, but not a solid code since $u_1$ and $u_2$ have a
	non-empty proper overlap.
	Then we have
	$
	\delta'(Q \times \{aa\}, u_2) = \delta(Q, a_1) \times \{a\}  	
	$
	and so $\delta(Q, a_1) \times \{a\} \subseteq \delta'(Q \times \pref{C} \setminus C, u_2)$.
	However, there are codes that are not solid codes for which it does hold.
	For example, for $C = \{ a_1, \ldots, a_{k-1}, a_k^m a_{k-1}\}$ for distinct $a_1, \ldots, a_k \in \Sigma$
	and $m \ge 1$.
\end{remark}

\begin{remark}
	Solid codes and the construction from the proof of  \cref{thm:constr_sync_hom_strongly_self_sync}
	can also be used to give an alternative proof of \cref{prelim:prop:sync_NL_complete}
	by first giving an automaton derived from the digraph
    by labeling the edges with letters and using an additional letter
    to realize the transitions going back to the start state.
    Then using a binary solid code
    and the construction from the proof of \cref{thm:constr_sync_hom_strongly_self_sync}
    we can reduce the alphabet to a binary alphabet.
    However, note that the words $bb,a,ba$ actually used in the proof of \cref{prelim:prop:sync_NL_complete}
    do not form a solid code.
\end{remark}

\begin{remark}
	We actually have $\Syn(\mathcal A') \cap \varphi(\Sigma^*) = \varphi(\Syn(\mathcal A))$
	in the construction from the proof of \cref{thm:constr_sync_hom_strongly_self_sync}.
\end{remark}

Coding morphism of solid codes are not the only morphisms
that preserve the complexity, as shown in the following remark.

\begin{remark}
	Consider the coding morphism $\varphi \colon \{a,b\}^* \to \{c,d, e\}^*$
	given by $\varphi(a) = cd$ and $\varphi(b) = ec$.
	Now $\varphi(\{a,b\}) = \{cd, ec\}$ is not a solid code.
	For $\semiautomat{A}{Q}{\{a,b\}}{\delta}$
	let 
	$\semiautomat{A'}{Q'}{\{c,d\}}{\delta'}$
	with $Q' = Q \times \{\varepsilon, c, e\}$
	where, for $q \in Q$, $u \in \{\varepsilon,c,d\}$, and $x \in \Sigma$,
	\[
	\delta'((q, u), x)
	= 
	\begin{cases}
		(\delta(q, a), \varepsilon)  & \text{if } u \in \{\varepsilon,c\} \text{ and } x = d, \\
		(\delta(q, b), \varepsilon)  & \text{if } u = e \text{ and } x = c, \\
		(q,x)                        & \text{if } u \in \{\varepsilon,c,e\} \text{ and } x \in \{c,e\} \text{ and } ux \ne ec \\
		%(q,e)                       & \text{if } u \in \{\varepsilon,e\} \text{ and } x = e,
		(q,u)                        & \text{if } u = e \text{ and } x = d.
	\end{cases}
	\]
	Then we have $\delta'((q, u), \varphi(v)) = (\delta(q, v), \varepsilon)$
	for $v \in \{a,b\}^+$ with $u \in \{\varepsilon, c, e\}$
	and $q \in Q$.
	This yields 
	\begin{align*}
		\delta'(Q', \varphi(a)) = \delta(Q, a) \times \{\varepsilon\}, \\
		\delta'(Q', \varphi(b)) = \delta(Q, b) \times \{\varepsilon\},	
	\end{align*}
	and	that for $|Q| \ge 2$ and regular $L \subseteq \{a,b\}^*$,
	the automaton~$\mathcal A$ has a synchronizing word in~$L$
	if and only if~$\mathcal A'$ has a synchronizing word in~$\varphi(L)$.
\end{remark}